%! Solving Exponential Equations with a Common Base — Unit 6, Lesson 6.3 — Warm-Up (Answer Key)
\documentclass[10pt]{article}
\usepackage{algebra2-article}
\usepackage{algebra2-key}   % pulls in -boxes; do NOT also load -boxes
\newcommand{\TallMath}[1]{$\displaystyle #1\rule[-1.4em]{0pt}{3.2em}$}

\begin{document}

\pageheader{Unit 6, Lesson 6.3}{Warm-Up --- Answer Key}
\namedateperiod

\vspace{0.10in}

\begin{notesbox}{Getting Ready: Skills We'll Use Today}
\small Today you \emph{solve} an exponential equation for the first time. Nothing below is new: two of
these items are Unit 5 exponent work, and the third is Lesson 6.0.

\vspace{5pt}
\textbf{1. The same number in a different outfit.} Fill in each missing exponent.

\vspace{4pt}
\hspace*{1.2em} $4 = 2^{\,\blacksquare}$, \; $\blacksquare=$ \ans{2} \qquad
$32 = 2^{\,\blacksquare}$, \; $\blacksquare=$ \ans{5} \qquad
$\dfrac{1}{8} = 2^{\,\blacksquare}$, \; $\blacksquare=$ \ans{$-3$} \qquad
$\sqrt{2} = 2^{\,\blacksquare}$, \; $\blacksquare=$ \ans{$\frac12$}

\vspace{5pt}
\hspace*{1.2em} $81 = 3^{\,\blacksquare}$, \; $\blacksquare=$ \ans{4} \qquad
$\dfrac{1}{27} = 3^{\,\blacksquare}$, \; $\blacksquare=$ \ans{$-3$} \qquad
$125 = 5^{\,\blacksquare}$, \; $\blacksquare=$ \ans{3} \qquad
$\sqrt[3]{7} = 7^{\,\blacksquare}$, \; $\blacksquare=$ \ans{$\frac13$}

\vspace{0.32in}

\textbf{2. Two rules that look alike and are not.} Evaluate each one \emph{as a number}, then say what
happened to the exponents.

\vspace{4pt}
\hspace*{1.2em} $2^{3}\cdot 2^{4} =$ \ans{128} $\;=\; 2^{\,\blacksquare}$ where $\blacksquare=$
\ans{7} \qquad
$\left(2^{3}\right)^{4} =$ \ans{4096} $\;=\; 2^{\,\blacksquare}$ where $\blacksquare=$
\ans{12}

\vspace{4pt}
\hspace*{1.2em} Multiplying powers of the same base \ans{adds} the exponents; raising a power to a
power \ans{multiplies} them.

\vspace{0.38in}

\textbf{3. A question about shape} (Lesson 6.0). Complete the table for $y=2^{x}$.

\vspace{3pt}
\begin{center}
\renewcommand{\arraystretch}{1.35}
\begin{tabular}{|c|c|c|c|c|c|}
\hline
$x$ & $-1$ & $0$ & $1$ & $2$ & $3$ \\ \hline
$y=2^{x}$ & \ans{$\frac12$} & \ans{1} & \ans{2} & \ans{4} & \ans{8} \\ \hline
\end{tabular}
\end{center}

\vspace{2pt}
\hspace*{1.2em} Is there any output this function produces \emph{twice} --- two different inputs giving
the same $y$? \ans{no}

\vspace{4pt}
\hspace*{1.2em} What was it about the \emph{shape} of an exponential graph in Lesson 6.0 that settles
that question? \ansline{It is always increasing --- no turning points, no maximum or minimum --- so it never comes back to an output it has already used.}

\end{notesbox}

\begin{teachernote}
\textbf{Item 2 is the signature error of the day, defused numerically before it can happen.} Students
who compute $128$ and $4096$ cannot then claim that $\left(2^{3}\right)^{x}$ is $2^{3+x}$. Write both
results on the board and leave them up --- Section 2 of the notes points straight back at them, and
every wrong answer in tonight's homework traces to this line.

\textbf{Item 3 is the entire justification for today's method, and it is worth two extra minutes.} The
answer ``no'' is not a fact about the table; it is a fact about the shape. Push until someone says
\emph{no turning points} or \emph{always increasing} (Lesson 6.0). Then say the consequence out loud
without proving it yet: if an exponential never repeats an output, then $2^{m}=2^{n}$ can only happen
when $m=n$ --- which is the permission slip for everything that follows. The Hook earns it with a
picture.

\textbf{Item 1 should be quick, and the last entry in each row is the one to check.} $\sqrt{2}=2^{1/2}$
and $\sqrt[3]{7}=7^{1/3}$ are Lesson 5.1, and they are the reason Tier E can ask for
$4^{x}=\sqrt{8}$. If several students stall on $\frac18=2^{-3}$, do that one at the board --- the
negative exponent shows up again immediately in $3^{x}=\frac{1}{27}$.
\end{teachernote}

\end{document}
